\metatitle{Weighted-bond symmetric functions}
\metadescription{Möbius and chromatic symmetric functions obtained from
multiweighted bond posets, with examples, recurrences, positivity results,
and open problems.}
\metakeywords{weighted bond poset, Möbius symmetric function, chromatic
symmetric function, chordal graph, parking function, Schur log-concavity}


\section[weightedBondSymmetric]{Weighted-bond symmetric functions}

\name{Rafael S. González D'León} and \name{Michelle L. Wachs} associate two
symmetric functions with the multiweighted bond posets of a graph
\cite{DLeonWachs2026x}.  The first refines the Möbius invariant of the bond
lattice; the second refines the chromatic polynomial.  The latter is different
from \name{Richard Stanley}'s
\hyperref[chromaticSymmetricFunction]{chromatic symmetric function}.


\subsection[weightedBondPoset]{Multiweighted bond posets}

Let \(G=(V,E)\).  Its \defin{bond lattice}\label{bondLattice} \(\Pi_G\) is the
subposet of the partition lattice consisting of set partitions whose blocks
induce connected subgraphs of \(G\).

For \(r\in\setP\), an \defin{\(r\)-weighted partition} has the form

\[
  \boldsymbol{\pi}
  =\{B_1^{\nu_1},\dotsc,B_k^{\nu_k}\},
\]

where \(\{B_1,\dotsc,B_k\}\) is a set partition and each \(\nu_i\) is a
\hyperref[weakCompositionDefinition]{weak \(r\)-composition} of
\(\lvert B_i\rvert-1\).  Refinement is allowed to merge blocks; when several
blocks are merged, the new weight must dominate the sum of their old weights.
Equivalently, a cover merges two blocks and adds one standard basis vector to
the sum of their weights.

The \defin{\(r\)-weighted bond poset} \(\Pi_G^r\) is the induced subposet on
weighted partitions whose underlying partition lies in \(\Pi_G\).  We write
\(\Pi_G^\infty\) when the weights are weak compositions with arbitrarily many
parts.  If \(G\) has \(n\) vertices and \(k\) connected components, then
\(\Pi_G^r\) is pure of length \(n-k\).


\subsection[mobiusSymmetricFunction]{The Möbius symmetric function}

\begin{polydata}{mobiusSymmetricFunction}
  Name     & Möbius symmetric functions of graphs \\
  Space    & Sym \\
  Basis    & False \\
  Rating   & 2 \\
  Bib      & DLeonWachs2026x \\
  Year     & 2026 \\
  Symbol   & \(M_G(\xvec)\) \\
  Category & Other \\
\end{polydata}

For a maximal element \(\boldsymbol{\pi}\) of \(\Pi_G^\infty\), let
\(w(\boldsymbol{\pi})\) be the sum of its block weights.  The
\defin{Möbius symmetric function of \(G\)} is

\[
  M_G(\xvec)
  \coloneqq
  \sum_{\boldsymbol{\pi}\in\operatorname{Max}(\Pi_G^\infty)}
  \mu_{\Pi_G^\infty}(\widehat{0},\boldsymbol{\pi})
  \xvec^{w(\boldsymbol{\pi})}.
\]

It is a homogeneous symmetric function of degree \(n-k\).  If \(G\) is
connected, symmetry gives the equivalent monomial expansion

\[
  M_G(\xvec)
  =
  \sum_{\lambda\vdash n-1}
  \mu_{\Pi_G^\infty}(\widehat{0},V^\lambda)
  \monomial_\lambda(\xvec).
\]

The specializations \(M_G(1,0,0,\dotsc)\) and \(M_G(1,t,0,\dotsc)\) are,
respectively, the Möbius invariant of \(\Pi_G\) and the Möbius polynomial of
the ordinary weighted bond poset.

% Related Rust: rust/sym-poly/sym/examples/weighted_bond_symmetric_site_example.rs
\begin{example}
Let \(P_3\) be the path on three vertices and \(K_3\) the complete graph.  The
recurrence below, computed exactly in the related Rust example, gives

\[
  M_{P_3}=\monomial_2+3\monomial_{11}
  =\elementaryE_{11}+\elementaryE_2,
\]

and

\[
  M_{K_3}=2\monomial_2+5\monomial_{11}
  =2\elementaryE_{11}+\elementaryE_2.
\]
\end{example}

The construction is multiplicative over connected components.  If \(G\) is
connected and has more than one vertex, then

\begin{align*}
  M_G
  &=-\sum_{\pi\in\Pi_G\setminus\{\widehat{1}\}}
    \completeH_{\lvert\pi\rvert-1}
    \prod_{B\in\pi}M_{G|B},\\
  M_G
  &=-\sum_{\pi\in\Pi_G\setminus\{\widehat{0}\}}
    M_{G/\pi}
    \prod_{B\in\pi}\completeH_{\lvert B\rvert-1}.
\end{align*}

For the path \(P_n\), this family recovers the
\hyperref[parking]{parking function symmetric function}:

\[
  (-1)^{n-1}M_{P_n}=\omega\parking_{n-1}
  =\sum_{\pi\in NC_{n-1}}\elementaryE_{\lambda(\pi)}.
\]

For a chordal graph with \(n\) vertices and \(k\) connected components,
\((-1)^{n-k}M_G\) is elementary-positive.  The corresponding statement for
arbitrary graphs remains open.

\begin{conjecture}[González D'León--Wachs]
For every graph \(G\) with \(n\) vertices and \(k\) connected components,
\((-1)^{n-k}M_G\) is
\hyperref[elementaryE]{elementary-positive}.
\end{conjecture}


\subsection[weightedBondChromaticSymmetric]{A chromatic symmetric analogue}

\begin{polydata}{weightedBondChromaticSymmetric}
  Name     & Weighted-bond chromatic symmetric functions \\
  Space    & Sym \\
  Basis    & False \\
  Rating   & 2 \\
  Bib      & DLeonWachs2026x \\
  Year     & 2026 \\
  Symbol   & \(\Psi_G(\xvec)\) \\
  Category & Other \\
  Contains & mobiusSymmetricFunction | DLeonWachs2026x | scope=top homogeneous component \\
\end{polydata}

Define

\[
  \Psi_G(\xvec)
  \coloneqq
  \sum_{\boldsymbol{\pi}\in\Pi_G^\infty}
  \mu_{\Pi_G^\infty}(\widehat{0},\boldsymbol{\pi})
  \xvec^{w(\boldsymbol{\pi})}.
\]

Equivalently,

\[
  \Psi_G(\xvec)=\sum_{\pi\in\Pi_G}M_{G|_\pi}(\xvec),
\]

where \(G|_\pi\) is the spanning subgraph whose connected components are the
induced graphs \(G|B\), for \(B\in\pi\).  Its highest homogeneous component is
\(M_G\), and its one-variable specialization reverses the chromatic
polynomial:

\[
  \Psi_G(t,0,0,\dotsc)=t^n\chi_G(t^{-1}).
\]

For the two three-vertex examples above,

\begin{align*}
  \Psi_{P_3}
    &=1-2\elementaryE_1+\elementaryE_{11}+\elementaryE_2,\\
  \Psi_{K_3}
    &=1-3\elementaryE_1+2\elementaryE_{11}+\elementaryE_2.
\end{align*}

In particular, their one-variable specializations are the reversals of
\(\chi_{P_3}(t)=t(t-1)^2\) and
\(\chi_{K_3}(t)=t(t-1)(t-2)\), respectively.

Write \(f|_d\) for the homogeneous component of degree \(d\).  The function
\(f\) is \defin{alternating elementary-positive} if
\((-1)^d f|_d\) is elementary-positive for every \(d\).  It is
\defin{Schur-log-concave} if

\[
  (f|_d)^2-(f|_{d-1})(f|_{d+1})
\]

is Schur-positive for every \(d\).  Chordal graphs give alternating
elementary-positive \(\Psi_G\); both of the following extensions are open.

\begin{conjecture}[González D'León--Wachs]
For every graph \(G\), the symmetric function \(\Psi_G\) is alternating
elementary-positive and Schur-log-concave.
\end{conjecture}

The Schur-log-concavity conjecture has been checked for connected graphs with
at most seven vertices.  For example, the degree-one inequality for \(P_3\)
is

\[
  (-2\elementaryE_1)^2
  -(\elementaryE_{11}+\elementaryE_2)
  =3\schurS_{11}+2\schurS_2.
\]

Finally, the first two equal-variable specializations for the complete graph
are

\[
  \Psi_{K_n}(t,0,\dotsc)=\prod_{i=1}^{n-1}(t-i),
  \qquad
  \Psi_{K_n}(t,t,0,\dotsc)=(t-n)^{n-1}.
\]

Finding comparably simple formulas after setting \(r>2\) variables equal to
\(t\) is another open problem.
